\section{Hoja de ruta de producto: Grace Blackwell, Vera Rubin y AI Factory}

\subsection{Señales de demanda detrás del cambio de hoja de ruta}

Jensen menciona en la entrevista que la velocidad del cambio de arquitectura de Grace Blackwell a Vera Rubin es alta porque la carga de trabajo está cambiando. Más allá del entrenamiento, la inferencia, la recuperación, la orquestación de múltiples agentes y el acceso a datos empresariales están modificando los cuellos de botella del sistema; la plataforma debe actualizarse conforme a la carga real, no seguir un ritmo fijo de lanzamientos.

\begin{importantbox}{La hoja de ruta no es un calendario, sino un bucle de retroalimentación}
La calidad de la hoja de ruta de producto depende de si la carga de los clientes, la investigación de frontera y la viabilidad de la cadena de suministro pueden retroalimentar las decisiones de arquitectura en tiempo real. Lo que Jensen describe con ``listening to whispers'' es precisamente ese bucle de retroalimentación.
\end{importantbox}

\subsection{El cambio de unidad en AI Factory}

Más adelante enfatiza que la ``unit of computing'' ha pasado de una sola máquina, luego a un clúster, y ahora a una ``AI factory''. Esto implica que las métricas de evaluación también deben cambiar: no basta con ver la velocidad de entrenamiento, también hay que ver la capacidad de producción de tokens, el consumo de energía por unidad, la automatización operativa y el ciclo de entrega. El centro de datos se concibe como una ``fábrica que produce inteligencia'', una perspectiva de industrialización y no de una sola máquina.

\begin{knowledgebox}{Métricas operativas de AI Factory}
Un sistema de métricas operables suele incluir: tokens per joule, rack utilization, MTTR, job completion SLA, unit economics (costo por millón de tokens) y la pérdida por cambio de versión de modelo. Aunque la entrevista no desarrolla cada punto, la narrativa de Jensen apunta claramente hacia estas dimensiones operativas.
\end{knowledgebox}

\subsection{Resumen de la sección}

La hoja de ruta de producto de NVIDIA ya no avanza en un solo eje de ``generaciones de chip'', sino de manera sincronizada bajo los objetivos multidimensionales de AI Factory. Lo que la nueva generación de sistemas debe responder no es ``qué tan rápido'', sino ``qué tan rápido y estable en producción real''.
